% Gliederung Stand 30.10.2026. ~14 Seiten.
% Je Unterkapitel unnummeriert: Versuchsaufbau -> Ergebnisse -> Diskussion. Erst Simulation, dann Real, gleiche Kennzahlen.
\chapter{Evaluation}
\label{chap:evaluation}
% Kurze Kapiteleinleitung vor 8.1 (Kap. 6/7 kurz zusammenfassen).

\section{Einheitliche Ausführung in Simulation und Realbetrieb}
\label{sec:eval_f1}
% F1, [A2]

\subsection{Vergleich von Code und Konfiguration}
\label{sec:eval_code}

\subsection{Realtests am Unitree G1}
\label{sec:eval_realtests}
% Inbetriebnahme, Gehen, Greifen, Übergaben.

\subsection{Anwendungsfall Werbevideo}
\label{sec:eval_video}
% Kompakt: genutzte Funktionen, was ohne Anpassung lief.

\section{Abbildungsgüte der Simulation}
\label{sec:eval_f2}
% F2, [A3, A4]

\subsection{Echtzeitfähigkeit}
\label{sec:eval_echtzeit}
% Regeltakte, Echtzeitfaktor auf Referenz-PC.

\subsection{Fortbewegung}
\label{sec:eval_fortbewegung}
% Geschwindigkeitsfolge, Stöße, Stand-Übergänge, Gangbild der Policies.

\subsection{Manipulation}
\label{sec:eval_manipulation}
% Planungserfolg und -zeit, IK-Restfehler, Jog-Genauigkeit Sim vs. Real.

\section{Sichere Bedienung ohne Vorerfahrung}
\label{sec:eval_f3}
% F3, [A5, A6]

\subsection{Sicherheitsfunktionen}
\label{sec:eval_sicherheit}
% Automatisierte Tests, Limits, Notaus-Reaktion.

\subsection{Bedienaufwand}
\label{sec:eval_bedienaufwand}
% Schritte und Zeit vom frischen PC bis zur ersten Bewegung.

\section{Erweiterbarkeit}
\label{sec:eval_f4}
% F4, [A1, A7]

\subsection{Anbindung eines externen Projekts}
\label{sec:eval_extern}

\subsection{Aufgabenunabhängigkeit}
\label{sec:eval_aufgaben}
% Neue Szene und neuer Ablauf ohne Codeänderung.

\section{Diskussion}
\label{sec:diskussion}

\subsection{Abgleich mit den Anforderungen}
\label{sec:abgleich}
% Tabelle A1–A7: erfüllt / teilweise / nicht erfüllt.

\subsection{Einordnung in den Stand der Technik}
\label{sec:einordnung}

\subsection{Limitierungen}
\label{sec:limitierungen}
% Kein Policy-Transfer, keine Real-Navigation, Sim-Kenngrößen nur teilweise real validiert.
